\section{Exploratory Data Analysis}

\subsection{Dataset overview and data model}

The dataset comprises \textbf{nine} related CSV tables totalling roughly 1.45
million rows. \code{orders} is the spine, at 99{,}441 rows -- one per order --
and the other eight attach to it directly or through \code{order\_items}.
Figure~\ref{fig:datamodel} shows the recovered structure and
Table~\ref{tab:dictionary} gives the full data dictionary.

Three tables are \textbf{one-to-many} against orders, and joining them naively is
the first trap the data sets: a direct orders--items--payments merge yields
118{,}437 rows rather than 99{,}441, a 19\% inflation that silently reweights
every mean towards large orders. We aggregate each child table to one row per
order first, then join, asserting the grain afterwards.

The second trap is identity. \code{customer\_id} is regenerated for every order,
so it is an order key, not a person: 99{,}441 distinct \code{customer\_id} values
map to only \textbf{96{,}096} distinct \code{customer\_unique\_id} values.
Grouping on the wrong one reports zero repeat buyers; the right one finds 2{,}997
people (3.1\%) who ordered again, one of them 17 times. We carry
\code{customer\_unique\_id} throughout and split on it
(Section~\ref{sec:candidates}).

\subsection{Data quality}

Table~\ref{tab:quality} records every issue and our decision; the decisions live
in \code{src/data\_load.py}, so the code and this table cannot drift apart.

Missingness here is \textbf{structural rather than random}. 97.0\% of orders reach
status \emph{delivered}; the rest are shipped, cancelled, unavailable or still
processing, and 2.98\% therefore carry no delivery timestamp. That is not a defect
to impute -- the event has not happened -- so delivery targets are undefined there
and excluded from delivery models. Elsewhere the data is remarkably complete:
99.2\% of orders carry a review and product attributes are missing for under 2\%
of products.

Two decisions shape the window. The first four months are a pilot contributing 326
orders against a 2018 run rate above 6{,}000 per month, and the final weeks are
truncated mid-cycle, so we restrict the analysis to \textbf{2017-01-01 to
2018-08-31}, leaving 99{,}091 orders (Figure~\ref{fig:coverage}).

\subsection{Temporal patterns: growth, seasonality and drift}
\label{sec:temporal}

Monthly order volume roughly triples across the window, peaking at \textbf{Black
Friday, November 2017} with 7{,}288 orders, about 1.5 times a typical month
(Figure~\ref{fig:volume}a). Purchasing is weekday-heavy, peaks at 16:00, and is
22\% weekend (Figure~\ref{fig:rhythm}).

The most consequential temporal finding is not seasonality but \textbf{drift}.
Mean lead time falls from 16.9 days in February 2018 to 7.7 in August 2018 --
Olist materially improved fulfilment over twenty months
(Figure~\ref{fig:volume}b) -- and across the chronological split used later the
mean falls by \textbf{5.34 days} between training and test.

Lateness behaves differently: it is \textbf{volatile rather than trending},
swinging from 1.4\% (June 2018) to 21.4\% (March 2018). Those spikes align with
Black Friday congestion and the nationwide truckers' strike that paralysed
Brazilian road freight in May 2018 -- exogenous shocks no order-level feature
available at checkout can anticipate.

Together these facts set the evaluation strategy. Under a level shift this size,
$R^2$ against a baseline fitted on the past is negative even for a genuinely
useful model, so we report \textbf{mean absolute error} as the headline regression
metric \citep{hyndman2006accuracy}, and treat periodic refitting as a Phase 2
requirement rather than an optimisation \citep{gama2014drift}.

\subsection{Customers, sellers, products and geography}

The marketplace is geographically \textbf{lopsided}, and that asymmetry drives the
delivery problem: \textbf{59.7\%} of sellers sit in S\~{a}o Paulo state against
\textbf{42.0\%} of customers, only 35.7\% of orders stay within the customer's own
state, and the median customer--seller distance is \textbf{434\,km}
(Figure~\ref{fig:geography}). Distance is the strongest checkout-time correlate of
lead time ($r = 0.40$), ahead of the platform's own estimate ($r = 0.38$).

Demand is long-tailed across 74 categories, the top ten covering two-thirds of
orders. The typical order is small -- median item value R\$86.90 -- but
\textbf{freight is not a rounding error}, a median R\$17.17 or 22\% of item value.
Payment is characteristically Brazilian: 75.4\% credit card, 51.4\% in instalments
(Figure~\ref{fig:products}).

\subsection{Reviews, and the relationship that dominates}
\label{sec:reviews}

Review scores are severely imbalanced -- \textbf{77.1\%} are four or five stars --
exactly the imbalance the brief warns about. The interesting structure is what
produces the minority.

A late delivery raises the probability of a 1--2 star review from \textbf{9.2\%}
to \textbf{54.0\%}, a \textbf{5.9-fold} increase, monotone in how late the parcel
is (Figure~\ref{fig:reviews}). This is the clearest relationship in the dataset,
and it reproduces at order level the logistics--rating link
\citet{deshpande2023logistics} establish across 15 million orders -- where it
carries through to seller sales. It also has a methodological consequence: a model
predicting review score \emph{at checkout} cannot see the variable that matters
most.

\subsection{Candidate problem exploration}
\label{sec:candidates}

We examined three candidates. All probes use \code{scikit-learn} pipelines on a
chronological split (train to 2018-04-30, validation to 2018-06-30, test
thereafter) additionally \textbf{grouped on \code{customer\_unique\_id}}, removing
the 227 buyers who would otherwise appear on both sides.\footnote{\textbf{Note on splitting.} We split by time rather than by
stratified random sampling because the target drifts (Section~\ref{sec:temporal})
and a random split lets a model learn from the future. Chronology alone is
insufficient though: 2{,}997 people placed more than one order, so a temporal cut
can still put one of a person's orders in training and another in test. This is
the concrete form of the brief's \code{customer\_id} versus
\code{customer\_unique\_id} warning, and why we also report a test subset of
customers never seen in training.}

\paragraph{Leakage depends on when the model runs.}
Rather than one global rule we define a forbidden set per \textbf{prediction
moment}. \emph{At checkout} nothing about fulfilment or the customer's opinion
exists: delivery dates, review scores and even \code{order\_approved\_at} are
future. \emph{At delivery} the parcel has arrived and we are predicting a review
not yet written, so the delivery outcome is history and admissible -- the brief
itself lists ``delivery performance'' among the features for this problem. The
guard \code{assert\_no\_leakage(features, regime)} enforces the right set, and
every modelling cell must pass it.

The cost of getting this wrong is measurable: adding one forbidden column, the
realised delivery time, lifts the late-delivery model's ROC-AUC from \textbf{0.589
to 1.000} (Figure~\ref{fig:leakage}). That is not a result; it is the target
restated.

\paragraph{Candidate 1 --- delivery performance, framed twice (primary).}
One question answered as a regression on lead time and a classification of late
versus on-time, satisfying the brief's requirement for both task types from one
coherent problem. Stakeholder: operations.

The regression works. Gradient boosting reaches \textbf{MAE 3.89 days} against
6.65 for a constant fitted on the training period, a \textbf{41.5\%} reduction,
with ridge at 4.61 days in between. Plain least squares, the simpler starting
point, is unusable here -- the design matrix is rank-deficient, so its
coefficients explode -- and ridge at $\alpha = 1$ is the minimum regularisation
that repairs it. Notably, Olist's own delivery promise is a
\emph{worse} point forecast than the constant (MAE 9.79 days) -- expected rather
than anomalous, since a promised date is a commercial decision, padded because a
broken promise costs returns and retention \citep{cui2024promise}. A useful
feature, not a useful prediction.

The classification is weak, and Section~\ref{sec:temporal} explains why: balanced
accuracy 0.580, PR-AUC 0.121 against a 0.074 floor. Crucially, \textbf{always
predicting ``on time'' scores 0.926 accuracy} while catching not a single late
order -- the exact trap the brief describes.

\paragraph{Candidate 2 --- low review score, predicted at delivery (secondary).}
The strongest classifier we found: PR-AUC \textbf{0.402} against a 0.110 floor
(3.6$\times$), balanced accuracy \textbf{0.692} -- and identical on customers never
seen in training (PR-AUC 0.403), so it is not memorising buyers. The prediction
moment is what makes it work: the same target predicted \emph{at checkout} reaches
only PR-AUC 0.212. Stakeholder: customer experience, with a concrete action --
flag an at-risk order for outreach before the review lands.

\paragraph{Candidate 3 --- repeat purchase within 30 days (rejected).}
This fails twice over. The base rate is only 1.8\% -- 96.9\% of customers appear
exactly once -- and the probe reaches balanced accuracy 0.517, essentially chance.
Separately, \textbf{right-censoring destroys the evaluation}: at the more natural
90-day horizon \emph{no} labelled orders remain in the test window at all, a
structural limit of a 21-month extract that no better model repairs.
Table~\ref{tab:checklist} scores all three.

\subsection{Key insights}

\begin{enumerate}\itemsep0pt\parsep0pt\topsep2pt
\item Three of the nine tables are one-to-many against orders; a naive join
inflates the data 19\% and biases every average towards large orders.
\item \code{customer\_id} is an order key, not a person -- the difference between
reporting 0\% and 3.1\% repeat buyers.
\item Delivery performance \textbf{drifts}: lead time halves over the window, so
MAE rather than $R^2$ is the honest regression metric.
\item Lateness is \textbf{volatile, not predictable} (1.4--21.4\% monthly), which
bounds the late-delivery classifier near 1.6$\times$ the no-skill floor.
\item A late delivery raises the 1--2 star rate from 9.2\% to 54.0\%, the
strongest relationship in the data.
\item \textbf{The prediction moment is a design decision, not a detail}: it nearly
doubles achievable PR-AUC one way, and inflates ROC-AUC to a meaningless 1.000 the
other.
\end{enumerate}

We expect a tuned Candidate 2 classifier to reach balanced accuracy in the high
60s, and Candidate 1's regression MAE of 3.5--4 days. The largest driver of both,
month-to-month carrier conditions, is not observable at prediction time; a
late-delivery classifier reporting 95\% accuracy would be the majority-class rule
or the leaked variant above. The dashboard and repository are linked on the cover
page; both read through the same \code{src/data\_load.py} entry point as this
report, and \code{src/report\_stats.py} regenerates every number quoted here.
